%% kapa: exam=1 round=25 year=2014 period=1 form=A subject=economics subject_ko=경제원론 seq=17 number=57
%% source: Q-Net 감정평가사 기출문제 (cst00309) artlSeq=5101207 file=extracted/1008/5101207/2014년 제25회 감정평가사 1차시험 기출문제/1교시(A형).hwp
%% conversion: kordoc 4.24.51 -> tools/gen1.py (auto); review=proofread-sonnet; flags=none
\begin{question}{57}
이윤을 극대화하는 기업의 생산함수가 $Q=2L^{ 0.5 } K^{ 0.5 }$이고 단위당 노동($L$)비용은 2, 자본($K$)비용은 1이다. 이 기업의 총비용이 100이고 제품의 시장가격이 10인 경우 다음 설명 중 옳지 않은 것은? (단, 제품시장과 생산요소시장은 완전경쟁적이다.)
\begin{choices}
  \item 노동을 30단위 사용해야 한다.
  \item 한계기술대체율($MRTS_{ LK }$)의 크기는 2이다.
  \item 이윤이 극대화되는 산출량은 50단위가 넘는다.
  \item 최대한 얻을 수 있는 이윤은 50을 넘는다.
  \item 자본을 50단위 사용해야 한다.
\end{choices}
\end{question}
\answer{1}
